\documentclass[11pt,a4paper]{article}
\usepackage[utf8]{inputenc}
\usepackage[T1]{fontenc}
\usepackage[spanish]{babel}
\usepackage[margin=2.5cm]{geometry}
\usepackage{booktabs}
\usepackage{array}
\usepackage{tabularx}
\usepackage{listings}
\usepackage{xcolor}
\usepackage{enumitem}
\usepackage{parskip}
\usepackage{fancyhdr}
\usepackage{amsmath}
\usepackage{amssymb}

\definecolor{codegray}{RGB}{245,245,245}
\definecolor{codeblue}{RGB}{0,60,180}
\definecolor{codegreen}{RGB}{0,120,0}

\lstdefinestyle{log}{
  basicstyle=\ttfamily\small,
  backgroundcolor=\color{codegray},
  frame=single,
  framerule=0.4pt,
  xleftmargin=8pt,
  numbers=none,
  breaklines=true
}

\pagestyle{fancy}
\fancyhf{}
\rhead{\textcolor{gray}{\small TD7: Ingeniería de Datos — UTDT}}
\lhead{\textcolor{gray}{\small Segundo Parcial Modelo 5}}
\rfoot{\thepage}
\renewcommand{\headrulewidth}{0.4pt}

\setlist[enumerate,1]{label=\textbf{\alph*)}}

\begin{document}

\begin{center}
  {\LARGE\bfseries TD7 — Ingeniería de Datos}\\[6pt]
  {\large Universidad Torcuato Di Tella}\\[4pt]
  {\Large\bfseries Segundo Parcial — Modelo 5}\\[4pt]
  {\normalsize Duración: 2 horas \quad|\quad Puntaje total: 100 puntos}
\end{center}
\noindent\rule{\linewidth}{0.8pt}
\vspace{2pt}

\noindent\textbf{Instrucciones:} Responda todos los ejercicios. Justifique cada respuesta; las afirmaciones sin justificación no suman puntaje. No está permitido el uso de material de consulta ni dispositivos electrónicos.

\vspace{6pt}
\noindent\rule{\linewidth}{0.4pt}

% ─────────────────────────────────────────────────────────────
\section*{Ejercicio 1 — Arquitectura Lambda y Master Dataset \hfill\normalfont(28 puntos)}

\begin{enumerate}
  \item (6 pts) En la arquitectura Lambda se trabaja con dos ecuaciones:
  \[
    \texttt{batch view = function(all data)} \qquad
    \texttt{realtime view = function(realtime view, new data)}
  \]
  Explique qué representa cada una y por qué la segunda es \textbf{incremental} mientras que la primera es de \textbf{alta latencia}. ¿Por qué la speed layer puede \textbf{descartar} datos cada vez que el batch layer termina un ciclo?

  \vspace{4.5cm}

  \item (7 pts) El \textbf{master dataset} es ``el núcleo y fuente de la verdad'' de la arquitectura. Explique por qué es la única parte que debe cuidarse especialmente y cómo permite recuperar el sistema completo si se pierden todas las vistas de la serving y la speed layer.

  \vspace{4cm}

  \item (9 pts) Enuncie y explique las tres propiedades (no funcionales) de los datos del master dataset: \textbf{crudeza (rawness)}, \textbf{inmutabilidad} y \textbf{perpetuidad}. ¿Por qué la inmutabilidad ``no tendría sentido sin la perpetuidad''? ¿Cómo aporta la inmutabilidad a la \textbf{tolerancia a fallas humanas}?

  \vspace{5.5cm}

  \item (6 pts) El patrón de acceso del master dataset es \textbf{write once, read many}. A partir de él, justifique por qué un \textbf{file system distribuido} se adapta mejor como \textit{batch layer storage} que un \textbf{key/value store}. Mencione una limitación del file system distribuido.

  \vspace{4.5cm}
\end{enumerate}

\newpage

% ─────────────────────────────────────────────────────────────
\section*{Ejercicio 2 — Seguridad \hfill\normalfont(24 puntos)}

\begin{enumerate}
  \item (7 pts) Mencione las cuatro \textbf{estrategias/técnicas} de seguridad vistas (autenticación, control de acceso, encriptación, firewalls / sistemas de detección) y explique brevemente qué problema ataca cada una. Diferencie \textbf{autenticación} de \textbf{control de acceso}.

  \vspace{4.5cm}

  \item (9 pts) Explique qué es un ataque de \textbf{SQL Injection}, sobre qué debilidad de las aplicaciones se aprovecha y qué objetivos puede tener (obtener, alterar o destruir información). Dé un ejemplo concreto de una consulta vulnerable y de un input malicioso, y describa cómo se mitiga.

  \vspace{6cm}

  \item (8 pts) Explique las \textbf{medidas de control de inferencia}. Para una base de datos estadística de población que solo expone agregados (por grupo de edad, ingresos, etc.), describa cómo un usuario podría \textbf{inferir} información confidencial de un individuo a pesar de no tener acceso directo a sus registros, y por qué esto es un problema de seguridad.

  \vspace{5cm}
\end{enumerate}

\newpage

% ─────────────────────────────────────────────────────────────
\section*{Ejercicio 3 — Privacidad, Compliance y Anonimización \hfill\normalfont(24 puntos)}

\begin{enumerate}
  \item (7 pts) Defina \textbf{PII} y explique qué obligaciones impone a una organización el manejo de PII bajo normativas como GDPR. ¿Qué rol cumplen los \textbf{data catalogs} para rastrear ``dónde vive y fluye'' la PII en los pipelines?

  \vspace{4.5cm}

  \item (8 pts) Explique qué es un \textbf{data catalog} y cómo \textbf{data discovery} lo extiende hacia un entendimiento dinámico y en tiempo real del estado de los datos. Enumere los cuatro \textbf{niveles de madurez} de datos de una organización (de \textit{data reactive} a \textit{data fluent}).

  \vspace{5cm}

  \item (9 pts) Describa el \textbf{proceso de anonimización}: por qué debe partir de una evaluación previa de los datos y por qué deben documentarse los detalles del proceso. Explique las técnicas de \textbf{perturbación}, \textbf{$k$-anonimato} y \textbf{$\ell$-diversidad}, dando un ejemplo de aplicación de cada una. ¿Qué es un \textbf{cuasi-identificador}?

  \vspace{6cm}
\end{enumerate}

\newpage

% ─────────────────────────────────────────────────────────────
\section*{Ejercicio 4 — Calidad, Big Data y Password Hashing \hfill\normalfont(24 puntos)}

\begin{enumerate}
  \item (8 pts) La \textbf{calidad de datos} se define como la aptitud de un dato para ser usado por quien debe consumirlo. Explique cómo se \textbf{mide} la calidad mediante análisis \textbf{univariado} y \textbf{bivariado}, dando dos ejemplos concretos de cada uno. ¿Qué dimensión de calidad atacaría con cada tipo de análisis?

  \vspace{5cm}

  \item (8 pts) Las cuatro ``V'' del Big Data incluyen la \textbf{veracidad}. Explique a qué se refiere y mencione tres de las dimensiones de calidad asociadas a la veracidad. Luego, para un escenario de \textit{high-frequency trading}, indique cuál de las ``V'' es la más crítica y por qué.

  \vspace{5cm}

  \item (8 pts) Un sistema almacena las contraseñas de sus usuarios. Explique el flujo completo de \textbf{password hashing con salting}: qué se almacena al crear la cuenta, y qué pasos ocurren al validar la contraseña en cada login. Si dos usuarios eligen exactamente la misma contraseña, ¿sus hashes almacenados serán iguales o distintos? Justifique.

  \vspace{5cm}
\end{enumerate}

\end{document}
